\exnumber{00}
\extitle{UX VI - Ex 1 - Parcours utilisateur}
\exfiles{ex00/fichier Figma, nomme Name\_FirstName\_DeliverableName\_Date\_VersionNumber}
\exauthorize{Those named in the exercise}
\makeheaderfiles

\textbf{Contexte :}

Your company's client is a movie theaters chain called, Movie 4 all.

They want to launch a platform for people who want to watch movies but don't want to spend endless time selecting one.
You are replacing someone in the company who has already worked on the user research, the problem statement and the list of features they plan to develop. Your goal is to prototype and test these features.

\textbf{Instructions :}

Based on the 3 features that have come out of the ideation workshops led by your collegues, design 3 user flows (1 per feature).

Features:

\begin{itemize}
\item 1: Adding a filter based on their current mood
\item 2: Offer a selection of 3 movies to be watched
\item 3: Ask users to rate the movie after watching it for further recommendation
\end{itemize}

\textbf{Livrables :}

\begin{itemize}
\item 3 user flows on a board (digital or pictures of a paper board)
\item Déposez les documents sur Github, dans un zip nommé « ProjectTitle\_Name\_FirstName », et chaque livrable ainsi : « Name\_FirstName\_DeliverableName\_Date\_VersionNumber ». Exemple : « Doe\_John\_MidFiPrototype\_05242024\_V1 ».
\item Ajoutez un titre dans tous les documents écrits.
\item Soignez l'aspect global des livrables. Ils doivent être clairs, simples et lisibles. Vous pouvez regarder canva.com, sans surcharger le dessin de détails.
\end{itemize}

\textbf{Ressources :}

\begin{itemize}
\item User flow glossary \url{https://www.productplan.com/glossary/user-flow/}
\item Site flows vs User Flows \url{https://uxmovement.com/wireframes/site-flows-vs-user-flows-when-to-use-which/}
\item How to make a User Flow diagram \url{https://www.lucidchart.com/blog/how-to-make-a-user-flow-diagram}
\end{itemize}

\textbf{Outils recommandés :}

\begin{itemize}
\item Figma
\item Whimsical
\end{itemize}

